% Original synthetic teaching sample for TeXada; AGPL-3.0-only.
% No copied textbook content or real personal data.
\documentclass[12pt]{article}
\usepackage{amsmath,amssymb}
\usepackage[margin=2.2cm]{geometry}
\setlength{\parindent}{0pt}
\setlength{\parskip}{0.7em}
% Intentional syntax error: the final subscript group lacks its closing brace.
\title{Sequence Notes: Indices and Powers}
\author{LinguistsWantTech --- synthetic teaching sample}
\date{}
\begin{document}
\maketitle
\section*{1. A clean power expression}
A geometric sequence can be written as $b_n=b_1 q^{n-1}$.
\section*{2. The recurrence with a grouping error}
The next term is q times the current term. Only the final subscript
brace has deliberately been removed:
$\displaystyle b_{n+1}=q b_{n$.
\section*{3. Keep the index meaning}
The intended change restores the grouping of the current index.
Do not replace an index with an invented value merely because the
new formula parses.
\textbf{Teaching note.} This original damaged example tests a
specific missing brace, not guessing the meaning of an absent subscript.
Model success is not assumed; unresolved input must remain available.
\end{document}
